\documentclass[11pt]{article}
\usepackage[margin=1in]{geometry}
\usepackage{amsmath,amssymb,booktabs}
\title{A Two-Page Note on Option Pricing}
\author{Formula Snip example}
\date{}
\begin{document}
\maketitle

\section{The Black--Scholes model}

Let $S_t$ be the price of a stock that follows a geometric Brownian motion with drift $\mu$
and volatility $\sigma$. The price of a European call option with strike $K$ and maturity $T$ is
\begin{equation}
  C = S_0 N(d_1) - K e^{-rT} N(d_2),
  \label{eq:bs}
\end{equation}
where $r$ denotes the risk-free interest rate and $N$ is the standard normal distribution function.
The two arguments are
\begin{equation}
  d_1 = \frac{\ln(S_0/K) + (r + \sigma^2/2)T}{\sigma\sqrt{T}}
\end{equation}
and
\begin{equation}
  d_2 = d_1 - \sigma\sqrt{T}.
\end{equation}
Put--call parity links the call to the put price $P$ through $C - P = S_0 - K e^{-rT}$.

\section{Risk and return}

The capital asset pricing model states that
\begin{equation}
  \mathrm{E}[R_i] = r_f + \beta_i \left( \mathrm{E}[R_m] - r_f \right),
\end{equation}
where $r_f$ is the risk-free rate and $\beta_i$ is the sensitivity of asset $i$ to the market.
The Sharpe ratio of a portfolio with mean return $\mu_p$ and volatility $\sigma_p$ is
\begin{equation}
  SR = \frac{\mu_p - r_f}{\sigma_p}.
\end{equation}

\begin{table}[h]
\centering
\begin{tabular}{lrr}
\toprule
Asset & $\beta_i$ & Sharpe ratio \\
\midrule
Stock A & 0.80 & 0.42 \\
Stock B & 1.20 & 0.35 \\
Stock C & 1.50 & 0.28 \\
\bottomrule
\end{tabular}
\caption{Example betas and Sharpe ratios.}
\end{table}

\newpage
\section{Discounting}

A cash flow $X$ received at time $T$ is worth $X e^{-rT}$ today. With annual compounding at rate $y$
the present value of a bond paying coupon $c$ for $n$ years and face value $F$ is
\begin{equation}
  PV = \sum_{k=1}^{n} \frac{c}{(1+y)^k} + \frac{F}{(1+y)^n}.
\end{equation}
The continuously compounded return over a period is $\ln(S_T/S_0)$, and its expected value
under the risk-neutral measure is $(r - \sigma^2/2)T$.

\end{document}
